
\documentclass[../../main.tex]{subfiles}
\begin{document}

% 年度の大きな表示
\begin{center}
  {\fontsize{48pt}{58pt}\selectfont \textbf{1970}\normalsize\textbf{年度}}
\end{center}

\vspace{10mm}

% 本文

\vspace{15mm}

% 出題テーマと難易度の表
\noindent\textbf{出題テーマと難易度}

\vspace{5mm}

\begin{center}
  \shadowbox{%
    \begin{tabular}{|c|p{8cm}|c|}
      \hline
      \textbf{問題番号} & \textbf{テーマ} & \textbf{難易度} \\
      \hline
      第1問 & 二次関数の最小値と円 & \\
      \hline
      第2問 & 楕円と極方程式 & \\
      \hline
      第3問 & 複素数と集合の決定 & \\
      \hline
      第4問 & 三角関数と不等式（整数解） & \\
      \hline
      第5問 & 三角関数の不等式証明 & \\
      \hline
      第6問 & 微分方程式 & \\
      \hline
    \end{tabular}%
  }
\end{center}

\vspace{15mm}

% 解答
\noindent\textbf{解答}

\vspace{5mm}

\begin{center}
  \shadowbox{%
    \renewcommand{\arraystretch}{1.6}
    \begin{tabular}{|c|c|p{8cm}|}
      \hline
      \textbf{問題番号} & \textbf{小問} & \textbf{答え} \\
      \hline
      \multirow{2}{*}{第1問} & (1) & $\alpha=0,t=2$のとき最小値$1$ \\
      \cline{2-3}
                             & (2) & $(p,q)=(0,2)$ \\
      \hline
      \multirow{2}{*}{第2問} & (1) & $r=\dfrac{a(1-e^2)}{1+e\cos\theta}$ \\
      \cline{2-3}
                             & (2) & $\dfrac{2-e^2}{a^2(1-e^2)^2}$ \\
      \hline
      第3問 & - & $S=\{\pm1\},\ \{\pm1,\pm i\}$ \\
      \hline
      第4問 & - & $n=5$ \\
      \hline
      第5問 & - & 証明問題 \\
      \hline
      第6問 & - & $y=e^x+9e^{\frac{2}{3}x}$ \\
      \hline
    \end{tabular}%
    \renewcommand{\arraystretch}{1.0}
  }
\end{center}

\vspace{15mm}

% 解答時間
\noindent\textbf{試験形態}

\vspace{5mm}

\begin{center}
  \shadowbox{%
    \begin{tabular}{p{8cm}}
      （要確認）
    \end{tabular}%
  }
\end{center}

\end{document}
